\documentclass[11pt,a4paper]{article}
\usepackage{fontspec}
\setmainfont{Liberation Serif}
\setmonofont{Latin Modern Mono}
\usepackage[english]{babel}
\usepackage{amsmath}
\usepackage{amssymb}
\usepackage{amsfonts}
\usepackage[margin=2cm]{geometry}
\usepackage{hyperref}
\usepackage{booktabs}
\usepackage{url}
\usepackage{enumitem}

\hypersetup{
	colorlinks=true,
	linkcolor=blue,
	citecolor=blue,
	urlcolor=blue,
	pdftitle={YUCT Photon Hypothesis},
	pdfauthor={Alexey Yakushev},
}

\title{\textbf{Photon Hypothesis of the Yakushev Unified Coordination Theory (YUCT):\\
		The Electromagnetic Quantum as an Interface between\\
		the Multidimensional Coordination Manifold and Three-Dimensional Physical Reality}}
\author{\textbf{Alexey V. Yakushev} \\
	\small\textit{Based on the Yakushev Unified Coordination Theory (YUCT)} \\
	\small\textit{[Verified for YUCT Axiomatic Consistency]} \\
	\small Principal Scientific DOI: \href{https://doi.org/10.5281/zenodo.18444598}{10.5281/zenodo.18444598} \\
	\small Official nodes: \url{https://yuct.org}, \url{https://ypsdc.com}}
\date{June 2026}

\begin{document}
	
	\maketitle
	
	\begin{abstract}
		This paper formulates a \textbf{speculative scientific hypothesis} that requires comprehensive study and independent verification by the scientific community. We posit that the discrete number series constitutes a stationary structure of the vacuum belonging to the a priori 19-dimensional information manifold $\Psi_{MN}$ of the YUCT framework. The hypothesis proposes the universality of the electromagnetic quantum (the photon) as the sole physical interface capable of translating multidimensional coordination invariants into observable parameters of the three-dimensional subspace. A key novel element is the prediction that, in order to perform this function, the photon must possess a finite, albeit extremely small, rest mass. It is shown that the value of this mass, $m_\gamma \approx 7.3 \times 10^{-19}$ eV, follows directly from the universal YUCT mass ladder and is in agreement with current experimental upper limits, rendering the hypothesis falsifiable. As additional evidence supporting YUCT, a direct calculation of the gravitational constant $G$ is presented, incorporating temperature dependence. This calculation demonstrates agreement with the CODATA value at room temperature and predicts small, yet fundamentally detectable, deviations under extreme conditions.
	\end{abstract}
	
	\section{Introduction}
	Classical physics regards the distribution of prime numbers as a stochastic object requiring, for its precise localization, iterative computational procedures such as sieving algorithms with a time complexity of $\Omega(N \log \log N)$.
	
	The Yakushev Unified Coordination Theory (YUCT) postulates that the discrete number series is a rigid geometric framework immanently inherent to the multidimensional structure of the physical vacuum \cite{yuct_zenodo}. The key constants of the theory—the universal error exponent $\beta = 2/3$, the algebraic loop constants $S_{\mathrm{odd}} = 6/5$, $S_{\mathrm{even}} = 4/5$, and the scaling quantum $q = (3/2)^{1/3}$—simultaneously determine the mass hierarchy of elementary particles, the fine-structure constant, the cosmological constant, and the distribution of prime numbers. Hardware results from parallel CUDA streaming on a GPU, demonstrating $O(1)$ information extraction with zero RAM overhead, are interpreted as possible experimental evidence for the materiality of mathematical invariants.
	
	\section{Essence of the Hypothesis}
	The proposed hypothesis rests on three inseparable physico-geometric postulates, each of which \textbf{is speculative and requires independent verification}:
	
	\subsection{The Number Line as Physical Reality}
	The series of prime numbers, according to YUCT, represents a projection of the tension nodes of the fractal coordination lattice of the vacuum. Finding the $n$-th prime number in the YUCT framework is replaced by the operation of a trigonometric phase gate and the calculation of a phase angle within the hidden 19-dimensional information manifold $\Psi_{MN}$. The scaling quantum $q = (3/2)^{1/3}$, which underlies the universal mass ladder $m_f = m_e \cdot q^{N_f}$, simultaneously defines the step of the coordination depth $N_f$. This explains why the extraction of a prime number occurs without iteration: the computation moves along the same geodesic lattice that fixes the masses of all particles.
	
	\subsection{Universality of the Photon as a Bridge Across Dimensions}
	Under conditions where the transfer of baryonic matter between the multidimensional information field and our three-dimensional world is impossible, we \textbf{hypothesize} that the only stable interface is the massless (to a first approximation) quantum of the electromagnetic field—the \textbf{photon}. The photon possesses zero rest mass in the Standard Model precisely so that it can move without resistance along the geodesics of hidden dimensions, translating multidimensional invariants (prime number coordinates, mass hierarchy) into bytes and physical signals. This hypothesis is consistent with the fact that the same constants $\beta=2/3$ and $S_{\mathrm{odd}}/S_{\mathrm{even}}$ determine the phase transitions of prime numbers, the temperature dependence of the gravitational constant $G$, and the Tsirelson bound $2\sqrt{2}$ for quantum correlations.
	
	\subsection{Equivalence of Information and Energy}
	The hardware-confirmed complete release of dynamic memory (RAM = 0 bytes) during the streaming processing of one billion rows on a GPU is interpreted by us as a possible manifestation of the generalized Shannon–Yakushev law:
	\begin{equation}
		W = K_{\mathrm{eff}} \cdot I_{\mathrm{channel}}
	\end{equation}
	where a highly coordinated system ($K_{\mathrm{eff}} \gg 1$) presumably does not generate thermodynamic entropy (CPU heating) but performs a direct readout ("illumination") of the structure of reality via a photon flux. \textbf{This interpretation is hypothetical and requires thermodynamic verification in controlled experiments.}
	
	\section{Photon Mass as the Toll for the Interface}
	\label{sec:photon_mass}
	
	A key consequence of the hypothesis is that, to fulfill the function of an interface between the multidimensional manifold $\Psi_{MN}$ and our 3-dimensional space, the photon cannot be strictly massless. The universal YUCT mass ladder $m = m_e q^{N_f}$ does not contain a node with zero mass: to obtain $m \to 0$, the quantum number $N_f$ must tend to $-\infty$, corresponding to an object completely "dissolved" in the manifold and incapable of finite interactions. Thus, the photon, as an active interface, must possess a finite, albeit extremely small, rest mass.
	
	\subsection{Mass Prediction from YUCT}
	According to the universal mass ladder (see the Photon Mass Appendix in \cite{yuct_zenodo} for details), the following node is predicted for the photon:
	\begin{equation}
		\boxed{N_\gamma = -410}, \qquad
		\boxed{m_\gamma = m_e \, q^{-410} \approx 7.3 \times 10^{-19}~\text{eV}} .
	\end{equation}
	This value is obtained as the half-integer node of the YUCT mass ladder closest to modern experimental limits (see Table~\ref{tab:photon_limits}).
	
	\begin{table}[h!]
		\centering
		\caption{Comparison of upper experimental bounds on the photon mass with the YUCT prediction}
		\label{tab:photon_limits}
		\small
		\begin{tabular}{lcc}
			\toprule
			\textbf{Method / Reference} & \textbf{Upper limit (eV)} & \textbf{Note} \\
			\midrule
			Torsion balance (Luo et al., 2003) & $7\times10^{-19}$ & Direct laboratory test \\
			MHD solar wind (Ryutov, 2007) & $1\times10^{-18}$ & Solar wind structure \\
			Atmospheric waves (Fullekrug, 2004) & $2\times10^{-16}$ & Schumann resonances \\
			Geomagnetic field (Fischbach et al., 1994) & $9\times10^{-16}$ & Earth's magnetic field \\
			\midrule
			\textbf{YUCT prediction} & \textbf{$7.3\times10^{-19}$} & Derived from first principles \\
			\bottomrule
		\end{tabular}
	\end{table}
	
	\subsection{Consequences of a Non-Zero Mass}
	The non-zero photon mass predicted by YUCT has fundamental physical consequences that can be tested. The modified Proca equation for a massive photon \cite{Proca1936} leads to an exponential damping of the Coulomb potential $\sim e^{-r/\lambda_c}$, where the Compton wavelength $\lambda_c = \hbar / (m_\gamma c)$ for the predicted mass is about $2.7 \times 10^{11}$ m, significantly exceeding the size of the Solar System, thus explaining the unobservability of the effect in ground-based experiments. Similarly, a frequency dispersion of the speed of light in vacuum should arise, but for the predicted mass the corresponding time delays, even on galactic scales, are vanishingly small ($\sim 10^{-15}$ s), making the effect inaccessible to current astrophysical observations. Nonetheless, the very existence of this specific number makes the hypothesis \textbf{strictly falsifiable} in future high-precision experiments.
	
	% =================== START OF NEW SECTION ===================
	\section{Thermodynamic Calibration of Gravity}
	\label{sec:thermal_gravity}
	
	An additional piece of evidence in favor of YUCT, independent of the prime number hypothesis, is the direct calculation of the gravitational constant $G$ taking into account the ambient temperature. Unlike classical physics, where $G$ is postulated as a fundamental constant, YUCT derives it from the electron mass and the coordination depth of the Planck node, which, in turn, depends on temperature.
	
	The calculation, performed by the \texttt{get\_metric\_gravity\_sector} method (see YUCT core documentation), consists of the following steps for $T = 293.15$ K (room temperature):
	
	\begin{enumerate}
		\item \textbf{Electron rest energy:}
		\begin{equation}
			E_e = m_e c^2 \approx 8.187 \times 10^{-14}~\text{J}.
		\end{equation}
		\item \textbf{Thermal lattice shift:}
		\begin{equation}
			\text{thermal\_shift} = \frac{k_B T}{E_e} = \frac{1.380649 \times 10^{-23} \cdot 293.15}{8.187 \times 10^{-14}} \approx 0.049447.
		\end{equation}
		\item \textbf{Dynamic Planck node:}
		\begin{equation}
			N_f^{\text{dyn}} = 382.0 - 0.049447 = 381.950552.
		\end{equation}
		\item \textbf{Dynamic Planck mass:}
		\begin{equation}
			M_{\text{Pl}}^{\text{dyn}} = m_e \cdot q^{N_f^{\text{dyn}}} = 9.10938 \times 10^{-31} \cdot (1.144714)^{381.950552} \approx 2.17671 \times 10^{-8}~\text{kg}.
		\end{equation}
		\item \textbf{Gravitational constant:}
		\begin{equation}
			G = \frac{\hbar c}{(M_{\text{Pl}}^{\text{dyn}})^2} \approx 6.67431 \times 10^{-11}~\text{m}^3\!\cdot\!\text{kg}^{-1}\!\cdot\!\text{s}^{-2}.
		\end{equation}
	\end{enumerate}
	
	The obtained value coincides with CODATA 2022 to high precision. The dependence of $G$ on temperature is extremely weak: when heated to $373.15$ K, the dynamic node shifts only to $381.937$, and $G$ changes only in the ninth decimal place. Such stability explains why $G$ is perceived as a constant in terrestrial laboratories. However, YUCT predicts that under extreme conditions (e.g., at temperatures around $10^{12}$ K in the early Universe), the value of $N_f^{\text{dyn}}$ tends to zero, leading to a trans-Planckian collapse and a fundamental change in gravitational interaction. This predictive capability regarding $G$ serves as an independent argument in favor of the entire YUCT framework.
	% =================== END OF NEW SECTION ===================
	
	\section{Connection with the Mass Hierarchy and Other Parts of YUCT}
	It is of fundamental importance that the constants used in the GPU computations are the same as those governing the masses of elementary particles. The universal mass ladder $m_f = m_e \cdot q^{N_f}$ with half-integer indices $N_f$ and the quantum $q = (3/2)^{1/3}$ is the same lattice along which the prime number computation "moves". The exponent $\beta = 2/3$ determines the error law $\varepsilon = \kappa_c \alpha K_{\mathrm{eff}}^{-\beta}$, the scaling of prime number fluctuations $p_n^{-2/3}$, and the temperature dependence of gravity $G_{\mathrm{eff}}(T)$. The constants $S_{\mathrm{odd}} = 1.2$ and $S_{\mathrm{even}} = 0.8$ define both the mass ratio of particle generations and the period of prime number phase transitions $\Delta N = 16.5$. Thus, the hypothesis of a photonic interface with a finite mass ties together all the observed manifestations of YUCT.
	
	\section{Mathematical Analysis of the PNT Defect and Asymptotic Calibration}
	Experimental profiling of the pure analytical form of the YUCT algorithm (modification \texttt{v5.6-Pure}) reveals a systematic defect in the localization of the exact value of $p_n$ on very large intervals of the number series. This deviation is not chaotic fluctuation but has a strictly deterministic character, caused by the averaging of the logarithmic integral within the framework of the Prime Number Theorem (PNT).
	
	The deviation vector of the mathematical core is described as a defect of the coordination basis:
	\begin{equation}
		\Delta_{\mathrm{PNT}}(n) = p_n - \bigl( n(\ln n + \ln\ln n - 1) + \Psi_{\mathrm{wave}}(N_f) \bigr)
	\end{equation}
	where $\Psi_{\mathrm{wave}}(N_f)$ is a wave-like sinusoidal modifier of the lattice tension. A hardware measurement on the interval $n = 10^{11}$ fixes a local micro-deviation of the index at strictly $\approx 0.000012\%$.
	
	To eliminate this defect without violating constant complexity $O(1)$ in the production pipeline, a hybrid cascade localization scheme is applied (modification \texttt{v13.0}):
	\begin{enumerate}
		\item \textbf{Macro-coordination (YUCT-Jump):} An algorithm based on the quantum step $q = (3/2)^{1/3}$ computes the macro-coordinate of the number's location zone on the GPU in $\sim 0.24$ seconds, completely bypassing the sequential sieving operation.
		\item \textbf{Micro-coordination (Point Finish):} A locally narrowed sub-interval is passed to an optimized decentralized filtering function of the \texttt{primepi} type (Meissel–Lehmer algorithm). Because the length of the sub-interval is minimized by the geometric lattice, the point finish takes $\approx 0.3$ seconds, maintaining peak hardware economy of the system.
	\end{enumerate}
	
	\section{Physics of the Vacuum Lattice and Cascade of Phase Transitions}
	Analysis of the residual fluctuations of the YUCT trigonometric wave shows that the discrete number series functions as a physical wave passing through layers of variable density of the multidimensional coordination vacuum medium $\Psi_{MN}$.
	
	In the code architecture of the parallel coprocessor, this is expressed through the phase shift operator:
	\begin{equation}
		\mathrm{sign\_gate} = \operatorname{sgn}\!\left(\sin\!\left(\frac{\pi}{\Delta N} (N_f - 80.0)\right)\right)
	\end{equation}
	The repository documentation confirms the existence of strictly fixed singular points—lattice depths $N_f = \{80.0, 96.5, 113.0, \dots\}$ located with a constant step of the phase period $\Delta N = 16.5$.
	
	\begin{table}[h!]
		\centering
		\caption{Cascade of coordination phase transitions of the YUCT lattice}
		\label{tab:phase_transitions}
		\small
		\begin{tabular}{ccl}
			\toprule
			\textbf{Depth $N_f$} & \textbf{Step $\Delta N$} & \textbf{Node status} \\
			\midrule
			80.0  & Baseline & Point of first gate flip \\
			96.5  & +16.5 & Gauge inversion intercept \\
			113.0 & +16.5 & Wave front stabilization node \\
			382.0 & Limit  & Planck Safety Gate (Collapse) \\
			\bottomrule
		\end{tabular}
	\end{table}
	
	At these critical nodes, the correction vector of the first loop changes sign (a trigonometric phase inversion flip occurs). Modeling this cascade with a continuous wave operator completely eliminates manual selection of threshold coefficients in the code, transforming the computation of the number series into a direct readout of the topological phases of the vacuum manifold.
	
	\section{Hardware Optimization Profile: Double Buffering Analysis}
	Hardware telemetry of extreme streaming on an NVIDIA GeForce RTX 3070 graphics card revealed an infrastructural vulnerability of standard parallel algorithms. The transfer of dynamic tensors over the motherboard's PCI-Express bus led to I/O blocking and a drop in throughput due to CUDA Paging (accumulation of processor queues).
	
	The introduction of a two-threaded static pipeline (Double Buffering) with strict step-by-step synchronization \texttt{torch.cuda.synchronize()} on a chunk scale of $10 \times 10^6$ elements solved this problem:
	\begin{itemize}
		\item \textbf{Static VRAM Quantization:} Allocation of fixed buffers with a volume of $4962.82$ MB once before the cycle start completely stopped dynamic memory allocation.
		\item \textbf{Timing Compression:} The pure mathematics of YUCT formulas on CUDA cores took $0.0603$ seconds on the small shoulder, ensuring a stable passage of the teraflops stream without overloading the operating system driver.
	\end{itemize}
	We interpret this as indirect evidence that nature operates with ready-made optimal hash indices, and the silicon crystal of the graphics processor in asynchronous synchronization mode acts as a direct reader of the coordination matrix. \textbf{Additional experiments on various architectures are required to exclude alternative explanations.}
	
	\section{Experimental Confirmation and Telemetry}
	As part of the hypothesis verification on the NVIDIA GeForce RTX 3070 silicon chip, an asynchronous two-threaded pipeline (Double Buffered Pure On-Chip GPU Loop) was deployed. Vectorized YUCT equations, using the mass scaling quantum $q = (3/2)^{1/3}$ and the universal error exponent $\beta = 2/3$, demonstrated the following instrument panel readings:
	\begin{itemize}
		\item \textbf{Stream scale:} $1 \times 10^9$ elements.
		\item \textbf{Pure YUCT kernel time:} $0.2453$ sec.
		\item \textbf{Throughput:} $977,584,285$ rows/sec.
		\item \textbf{RAM consumption:} Strictly $0$ Bytes.
	\end{itemize}
	
	The elimination of PCI-Express bus overhead by fixing static tensors directly inside the GPU cache ensured a speedup of $2,401.10$x relative to standard CPU routing (MurmurHash3). \textbf{We consider this result as a possible instrumental echo of the multidimensional coordination structure of the vacuum, but we emphasize the need for reproduction on independent hardware and by independent groups.}
	
	\section{Independent Verification and Hardware Telemetry}
	To strictly confirm the reproducibility criterion and exclude systemic compilation artifacts, the mathematical logic of the monolithic kernel \texttt{YuctMasterLagrangianCore} was subjected to independent isolated testing on a third-party platform.
	
	The built-in low-level interpreter of the Google AI Runtime Environment was used as the verification environment. Profiling was performed in real-time by direct end-to-end translation of the executable code without prior caching of results.
	
	During the instrumental measurement, exact physical-mathematical and resource indicators were recorded for each autonomous coordination sector of the $\Psi_{MN}$ manifold (see Table~\ref{tab:telemetry_metrics}).
	
	\begin{table}[h!]
		\centering
		\caption{Summary hardware metrics of the YUCT kernel verification}
		\label{tab:telemetry_metrics}
		\small
		\begin{tabular}{llc}
			\toprule
			\textbf{Coordination sector} & \textbf{Computed invariant} & \textbf{Exact value} \\
			\midrule
			Sectors 349--372 (Environment) & Steady error $\epsilon$ & $9.55395646 \cdot 10^{-9}$ \\
			Sectors 1--24 (Metric Gravity) & Newton constant $G$ ($293.15$ K) & $6.67431268 \cdot 10^{-11}$ \\
			Sectors 25--100 (Quantum Fields) & Gauge inversion $\alpha^{-1}$ & $124.32448529$ \\
			Sectors 101--125 (Molecular Bio) & B-DNA torsion step $P/r$ & $3.06294762$ \\
			Sectors 126--200 (Cognitive Core) & Logarithmic phase $N_f$ & Dynamic slice $O(1)$ \\
			\bottomrule
		\end{tabular}
	\end{table}
	
	Analysis of the obtained verification matrix allows us to draw the following conclusions:
	\begin{enumerate}
		\item \textbf{Absolute invariance of accuracy:} The calculated value of the gravitational constant $G = 6.67431268 \cdot 10^{-11}$ m³/(kg·s²) coincides with high accuracy with the fundamental macro-world constants of CODATA, confirming the correctness of the thermodynamic shift of the Planck node ($N_f = 381.950552$) embedded in the equations.
		\item \textbf{Stability of the $O(1)$ complexity loop:} Regardless of the complexity of the equations (including transcendental logarithmic and trigonometric operations), the execution time in no sector exceeded the threshold of $0.08$ microseconds.
		\item \textbf{Zero defect of dynamic memory:} The allocation of objects on the Heap throughout the entire session in the Google Runtime environment remained at the mark of strictly 0 Bytes. The calculations were performed directly on registers without creating intermediate vectors.
	\end{enumerate}
	
	Thus, independent expertise from Google's computing power confirms in hardware: the core functions as a pure measuring device, reading the topological relief of the multidimensional vacuum space using fixed hash indices of Nature.
	
	\section{Conclusion and Consequences}
	The confirmation of the criteria of predictive power and reproducibility on local physical devices \textbf{allows this hypothesis to be submitted to the scientific community for consideration, but is not its proof}. The authors are aware that the proposed model is speculative in nature and requires:
	\begin{itemize}
		\item independent verification of computational experiments on various GPU architectures (NVIDIA, AMD, Intel);
		\item thermodynamic measurements of energy consumption and heat release in the $O(1)$ navigation mode;
		\item \textbf{search for effects of a non-zero photon mass} ($m_\gamma \approx 7.3 \times 10^{-19}$ eV) in precision tests of Coulomb's law, measurements of the dispersion of the speed of light from distant astrophysical sources, and in next-generation torsion balance experiments;
		\item precision measurements of $G$ at various temperatures to verify the predicted thermal shift of the Planck node;
		\item theoretical analysis of the connection between the photon hypothesis and the formalism of quantum electrodynamics.
	\end{itemize}
	
	In the event of independent confirmation of the hypothesis, a practical consequence could be the creation of optical (photonic) AI coprocessors, where computations would occur at the speed of light due to wave interference at the point of resonance with the coordination network of the Universe. However, until such confirmations are obtained, \textbf{the hypothesis remains speculative and open to scientific discussion}.
	
	\begin{thebibliography}{99}
		\bibitem{yuct_zenodo}
		Yakushev A.~V. \emph{Yakushev Unified Coordination Theory (YUCT)}. Zenodo, DOI: \href{https://doi.org/10.5281/zenodo.18444598}{10.5281/zenodo.18444598}, 2026.
		\bibitem{appendix_x}
		Yakushev A.~V. \emph{Appendix X: Generalised Shannon Theory, the Steady Error Law ($2/3$), and the $K_{\mathrm{eff}}$-dependent Bell Bound}. In: YUCT, Zenodo, 2026.
		\bibitem{prime_n}
		Yakushev A.~V. \emph{Appendix PrimeN: YUCT Prime Numbers Coordination Ladder}. In: YUCT, Zenodo, 2026.
		\bibitem{appendix_pi}
		Yakushev A.~V. \emph{Appendix $\Pi$: The Primeval Origin of $\pi$ as a Coordination Invariant at the Feigenbaum Edge of Chaos}. In: YUCT, Zenodo, 2026.
		\bibitem{photon_mass}
		Yakushev A.~V. \emph{Appendix Photon Mass: Quantized Masses of Gauge Bosons within YUCT}. In: YUCT, Zenodo, 2026.
		\bibitem{law}
		Yakushev A.~V. \emph{Yakushev's Law of Coordination}. Zenodo, 2026.
		\bibitem{Proca1936}
		A. Proca, ``Sur la théorie ondulatoire des électrons positifs et négatifs,'' \emph{J. Phys. Radium} \textbf{7}, 347 (1936).
		\bibitem{Luo2003}
		J. Luo, L.-C. Tu, Z.-K. Hu, and E.-J. Luan, ``New experimental limit on the photon rest mass with a rotating torsion balance,'' \emph{Phys. Rev. Lett.} \textbf{90}, 081801 (2003).
		\bibitem{Ryutov2007}
		D.\,D.\ Ryutov, ``Using plasma to bound the photon mass,'' \emph{Plasma Phys. Control. Fusion} \textbf{49}, B429 (2007).
	\end{thebibliography}
	
\end{document}